\documentclass[12pt, a4paper]{report}

% ─── Paquetes ───────────────────────────────────────────────────────────────
\usepackage[utf8]{inputenc}
\usepackage[T1]{fontenc}
\usepackage[spanish]{babel}
\usepackage{geometry}
\usepackage{graphicx}
\usepackage{hyperref}
\usepackage{listings}
\usepackage{xcolor}
\usepackage{tikz}
\usepackage{booktabs}
\usepackage{tabularx}
\usepackage{array}
\usepackage{enumitem}
\usepackage{fancyhdr}
\usepackage{titlesec}
\usepackage{mdframed}
\usepackage{float}
\usepackage{lmodern}
\usepackage{microtype}
\usepackage{parskip}

\usetikzlibrary{shapes.geometric, arrows.meta, positioning, fit, backgrounds, calc}

% ─── Geometría ──────────────────────────────────────────────────────────────
\geometry{top=2.5cm, bottom=2.5cm, left=3cm, right=2.5cm}

% ─── Colores ────────────────────────────────────────────────────────────────
\definecolor{gris}{RGB}{80,80,80}
\definecolor{grisCaja}{RGB}{245,245,245}
\definecolor{negro}{RGB}{20,20,20}
\definecolor{linea}{RGB}{180,180,180}
\definecolor{codigofondo}{RGB}{248,248,248}

% ─── Código ─────────────────────────────────────────────────────────────────
\lstset{
  backgroundcolor=\color{codigofondo},
  basicstyle=\ttfamily\small,
  breaklines=true,
  frame=single,
  framesep=4pt,
  rulecolor=\color{linea},
  xleftmargin=8pt,
  xrightmargin=8pt,
  showstringspaces=false,
  tabsize=2
}

% ─── Caja de nota ───────────────────────────────────────────────────────────
\newmdenv[
  linecolor=linea,
  backgroundcolor=grisCaja,
  linewidth=1pt,
  innerleftmargin=10pt,
  innerrightmargin=10pt,
  innertopmargin=8pt,
  innerbottommargin=8pt,
  skipabove=8pt,
  skipbelow=8pt
]{notacaja}

% ─── Encabezado ─────────────────────────────────────────────────────────────
\pagestyle{fancy}
\fancyhf{}
\fancyhead[L]{\small\textit{Spec-Driven Development con Spec Kit}}
\fancyhead[R]{\small\textit{Jhonatan Castro}}
\fancyfoot[C]{\thepage}
\renewcommand{\headrulewidth}{0.4pt}

% ─── Títulos ─────────────────────────────────────────────────────────────────
\titleformat{\chapter}[block]
  {\Large\bfseries}{\thechapter.}{10pt}{}
  [\vspace{4pt}\rule{\textwidth}{0.4pt}]
\titleformat{\section}{\large\bfseries}{\thesection}{8pt}{}
\titleformat{\subsection}{\normalsize\bfseries}{\thesubsection}{6pt}{}

\hypersetup{colorlinks=true,linkcolor=negro,urlcolor=gris,citecolor=gris}

% ════════════════════════════════════════════════════════════════════════════
\begin{document}

% ─── Portada ────────────────────────────────────────────────────────────────
\begin{titlepage}
  \centering
  \vspace*{3cm}
  {\Large\bfseries Spec-Driven Development}\\[0.4cm]
  {\large\bfseries Aprendizaje con GitHub Spec Kit}\\[2cm]
  \rule{\textwidth}{0.4pt}\\[0.5cm]
  {\normalsize
    Documentaci\'on del proceso de aprendizaje: metodolog\'ia SDD,\\
    herramienta Spec Kit y construcci\'on de un Sistema de Gesti\'on\\
    de Torneos como proyecto de pr\'actica.
  }\\[0.4cm]
  \rule{\textwidth}{0.4pt}
  \vspace{2cm}
  \begin{tabular}{ll}
    \textbf{Autor}   & Jhonatan Castro \\[4pt]
    \textbf{GitHub}  & JhonatanC4STRO \\[4pt]
    \textbf{Fecha}   & Junio 2026 \\[4pt]
    \textbf{Agente}  & Claude Code (Sonnet 4.6) \\
  \end{tabular}
  \vfill
  {\small Este documento fue generado a partir de una sesi\'on real de aprendizaje.}
\end{titlepage}

\tableofcontents
\newpage

% ════════════════════════════════════════════════════════════════════════════
\chapter{Introducci\'on}

\section{Contexto}

Este documento registra una sesi\'on de aprendizaje orientada a comprender
el \textbf{Spec-Driven Development (SDD)} como alternativa al denominado
\textit{vibe coding}: la pr\'actica de generar c\'odigo mediante prompts
improvisados a un agente de IA sin planificaci\'on previa.

El objetivo fue doble: entender la metodolog\'ia SDD a nivel conceptual,
y aplicarla de forma pr\'actica utilizando \textbf{GitHub Spec Kit} como
herramienta, \textbf{Claude Code} como agente de implementaci\'on, y un
\textit{Sistema de Gesti\'on de Torneos} como proyecto de pr\'actica.

\section{Objetivos de aprendizaje}

\begin{enumerate}[leftmargin=*, label=\arabic{*}.]
  \item Comprender la diferencia entre SDD y vibe coding.
  \item Distinguir entre la metodolog\'ia (SDD) y la herramienta (Spec Kit).
  \item Ejecutar el flujo completo: \texttt{constitution} $\to$
        \texttt{specify} $\to$ \texttt{plan} $\to$ \texttt{tasks}
        $\to$ \texttt{implement}.
  \item Evaluar alternativas a Spec Kit (GSD, BMAD, Kiro, OpenSpec).
\end{enumerate}

% ════════════════════════════════════════════════════════════════════════════
\chapter{Marco Conceptual}

\section{?`Qu\'e es el Spec-Driven Development?}

El \textbf{Spec-Driven Development (SDD)} es una metodolog\'ia de desarrollo
de software que trata las especificaciones como artefactos de primera clase.
En lugar de comenzar a generar c\'odigo inmediatamente, el desarrollador
define con precisi\'on \textit{qu\'e} debe construirse antes de que el agente
toque una sola l\'inea de c\'odigo.

\begin{notacaja}
  \textbf{Principio fundamental:} La especificaci\'on es la fuente de verdad.
  El c\'odigo es el \'ultimo paso, no el primero.
\end{notacaja}

\section{Vibe Coding vs.\ SDD}

\begin{table}[H]
\centering
\begin{tabularx}{\textwidth}{lXX}
\toprule
\textbf{Aspecto} & \textbf{Vibe Coding} & \textbf{SDD} \\
\midrule
Punto de partida & Prompt improvisado & Especificaci\'on detallada \\
Decisiones       & Las toma el agente & Las toma el desarrollador \\
Revisi\'on       & Al final (si ocurre) & En cada fase \\
Fuente de verdad & El c\'odigo generado & La especificaci\'on \\
Resultado        & Impredecible & Trazable \\
\bottomrule
\end{tabularx}
\caption{Comparaci\'on entre Vibe Coding y SDD}
\end{table}

\section{SDD vs.\ Spec Kit: una distinci\'on cr\'itica}

\begin{table}[H]
\centering
\begin{tabularx}{\textwidth}{lXX}
\toprule
& \textbf{SDD} & \textbf{Spec Kit} \\
\midrule
Naturaleza & Metodolog\'ia / filosof\'ia & Herramienta (CLI + templates) \\
Origen     & Concepto independiente & GitHub (open source, MIT) \\
Requisito  & Ninguno & Python 3.11+ \\
Analog\'ia & Scrum & Jira \\
\bottomrule
\end{tabularx}
\caption{SDD como metodolog\'ia vs.\ Spec Kit como herramienta}
\end{table}

Es posible practicar SDD sin Spec Kit, de la misma manera que es posible
gestionar un proyecto \'agil sin Jira. La herramienta implementa la
metodolog\'ia; no la sustituye.

% ════════════════════════════════════════════════════════════════════════════
\chapter{GitHub Spec Kit}

\section{Descripci\'on general}

GitHub Spec Kit es un conjunto de herramientas open source (MIT License)
que implementa el flujo SDD para agentes de IA. Alcanz\'o m\'as de 93.000
estrellas en GitHub hacia mayo de 2026, con soporte para m\'as de 30 agentes
incluyendo Claude Code, GitHub Copilot, Cursor y Gemini CLI.

\section{Arquitectura de archivos}

Al instalar Spec Kit en un proyecto, se genera la siguiente estructura:

\begin{lstlisting}
proyecto/
|-- constitution.md          <- ADN del proyecto (escrito por el dev)
|-- .claude/
|   `-- skills/
|       |-- speckit-constitution/
|       |-- speckit-specify/
|       |-- speckit-plan/
|       |-- speckit-tasks/
|       |-- speckit-implement/
|       `-- speckit-analyze/
|-- .specify/
|   |-- templates/
|   `-- memory/
`-- specs/
    |-- 001-nombre-feature/
    |   |-- spec.md
    |   |-- plan.md
    |   |-- tasks.md
    |   |-- research.md
    |   |-- data-model.md
    |   `-- contracts/
    `-- 002-nombre-feature/
\end{lstlisting}

\section{El flujo de comandos}

\begin{figure}[H]
\centering
\begin{tikzpicture}[
  nodo/.style={
    rectangle, rounded corners=4pt,
    minimum width=2.8cm, minimum height=1cm,
    draw=gris, fill=grisCaja,
    font=\small\ttfamily, align=center
  },
  flecha/.style={-Stealth, thick, color=gris}
]
\node[nodo] (c) {constitution};
\node[nodo, right=0.9cm of c] (s) {specify};
\node[nodo, right=0.9cm of s] (p) {plan};
\node[nodo, right=0.9cm of p] (t) {tasks};
\node[nodo, right=0.9cm of t] (i) {implement};

\draw[flecha] (c)--(s);
\draw[flecha] (s)--(p);
\draw[flecha] (p)--(t);
\draw[flecha] (t)--(i);

\node[below=0.45cm of c, font=\tiny, text=gris, align=center] {ADN del\\proyecto};
\node[below=0.45cm of s, font=\tiny, text=gris, align=center] {Define la\\feature};
\node[below=0.45cm of p, font=\tiny, text=gris, align=center] {Plan\\t\'ecnico};
\node[below=0.45cm of t, font=\tiny, text=gris, align=center] {Tareas\\peque\~nas};
\node[below=0.45cm of i, font=\tiny, text=gris, align=center] {Genera\\c\'odigo};
\end{tikzpicture}
\caption{Flujo de comandos en Spec Kit (ejecutados en Claude Code)}
\end{figure}

\section{El archivo \texttt{constitution.md}}

El \texttt{constitution.md} es el artefacto m\'as importante. Es el contrato
entre el desarrollador y el agente: define el stack, los principios
inamovibles, la estructura de carpetas, las convenciones de c\'odigo y
los patrones prohibidos. El agente lo respeta en cada decisi\'on.

\begin{lstlisting}[caption={Ejemplo de constitution.md}]
# Constitution -- Sistema de Gestion de Torneos

## Stack
- React (Vite) -- frontend
- Node.js + Express -- backend/API
- PostgreSQL + Prisma -- base de datos
- JWT manual con bcrypt -- autenticacion del admin
- Tailwind CSS -- estilos

## Principios inamovibles
- Nunca usar `any` en TypeScript
- Toda funcion debe tener su tipo de retorno explicito
- No instalar librerias sin aprobar primero
- Toda feature debe tener su spec antes de codear

## Estructura de carpetas
/client  -> React (frontend)
/server  -> Node.js + Express (backend)

## Patrones prohibidos
- No mezclar logica de negocio en componentes de UI
- No exponer contrasenas ni tokens en respuestas de la API
- No hardcodear IDs ni valores magicos
\end{lstlisting}

\section{Participaci\'on activa del desarrollador}

Un error frecuente es creer que el flujo SDD es autom\'atico. No lo es.
Cada comando requiere revisi\'on y aprobaci\'on antes de continuar al
siguiente:

\begin{table}[H]
\centering
\begin{tabularx}{\textwidth}{lX}
\toprule
\textbf{Comando} & \textbf{Qu\'e revisa el desarrollador} \\
\midrule
\texttt{/speckit-constitution} & Que el agente interpret\'o bien el ADN \\
\texttt{/speckit-specify}      & Que la spec captura todos los casos de uso \\
\texttt{/speckit-plan}         & Que las decisiones arquitect\'onicas son correctas \\
\texttt{/speckit-tasks}        & Que las tareas son peque\~nas y testeables \\
\texttt{/speckit-implement}    & Que el c\'odigo generado cumple la spec \\
\bottomrule
\end{tabularx}
\caption{Responsabilidad del desarrollador en cada fase}
\end{table}

% ════════════════════════════════════════════════════════════════════════════
\chapter{Proyecto de Pr\'actica: Sistema de Gesti\'on de Torneos}

\section{Descripci\'on}

Como proyecto de pr\'actica se eligi\'o un \textbf{Sistema de Gesti\'on de
Torneos} para dos juegos: FC~25 y Call of Duty Black Ops~2.

\begin{table}[H]
\centering
\begin{tabularx}{\textwidth}{lX}
\toprule
\textbf{Componente} & \textbf{Tecnolog\'ia} \\
\midrule
Frontend      & React + Vite + TypeScript + Tailwind CSS \\
Backend       & Node.js + Express + TypeScript \\
Base de datos & PostgreSQL + Prisma ORM \\
Autenticaci\'on & JWT manual + bcrypt \\
\bottomrule
\end{tabularx}
\caption{Stack tecnol\'ogico del proyecto de pr\'actica}
\end{table}

\section{Features especificadas}

El sistema se dividi\'o en 5 features procesadas de forma independiente:

\begin{enumerate}[leftmargin=*, label=\textbf{F\arabic{*}.}]
  \item \textbf{Inscripci\'on al torneo} --- Formulario p\'ublico con nombre
        completo, nickname, n\'umero de documento, n\'umero de ficha, nombre
        de formaci\'on y selecci\'on de juego. Validaci\'on de duplicados por
        nickname y documento dentro del mismo juego.
  \item \textbf{Login del administrador} --- Autenticaci\'on con email y
        contrase\~na. Token JWT con duraci\'on de 8 horas. Mensajes de error
        en caso de credenciales inv\'alidas.
  \item \textbf{Panel admin: gesti\'on de inscritos} --- Listado por juego,
        eliminaci\'on de jugadores, apertura y cierre de inscripciones.
  \item \textbf{Generaci\'on del bracket} --- FC~25 usa single-elimination;
        Call of Duty usa double-elimination. Asignaci\'on aleatoria de
        posiciones. Estructura inmutable una vez generada.
  \item \textbf{Registro de resultados} --- Ganador, puntos, empate y penales.
        Avance autom\'atico del ganador. Correcci\'on permitida solo si el
        partido siguiente a\'un no tiene resultado.
\end{enumerate}

\section{Flujo de actores}

\begin{figure}[H]
\centering
\begin{tikzpicture}[
  actor/.style={
    rectangle, rounded corners=3pt,
    minimum width=2.6cm, minimum height=0.8cm,
    draw=gris, fill=grisCaja,
    font=\small, align=center
  },
  accion/.style={
    rectangle,
    minimum width=4.5cm, minimum height=0.75cm,
    draw=gris, font=\small, align=center
  },
  flecha/.style={-Stealth, color=gris}
]

\node[actor] (jugador) {Jugador P\'ublico};
\node[accion, right=1.2cm of jugador]           (f1) {Formulario de inscripci\'on};
\node[accion, below=0.5cm of f1]                (f2) {Validaci\'on de datos};
\node[accion, below=0.5cm of f2]                (f3) {Guardar en PostgreSQL};

\node[actor, below=2.2cm of jugador] (admin) {Administrador};
\node[accion, right=1.2cm of admin]             (a1) {Login (JWT)};
\node[accion, below=0.5cm of a1]                (a2) {Ver inscritos por juego};
\node[accion, below=0.5cm of a2]                (a3) {Generar bracket};
\node[accion, below=0.5cm of a3]                (a4) {Registrar resultados};
\node[accion, below=0.5cm of a4]                (a5) {Avance autom\'atico};

\draw[flecha] (jugador)--(f1);
\draw[flecha] (f1)--(f2);
\draw[flecha] (f2)--(f3);
\draw[flecha] (admin)--(a1);
\draw[flecha] (a1)--(a2);
\draw[flecha] (a2)--(a3);
\draw[flecha] (a3)--(a4);
\draw[flecha] (a4)--(a5);
\end{tikzpicture}
\caption{Flujo de actores del Sistema de Gesti\'on de Torneos}
\end{figure}

\section{Formato de brackets por juego}

\begin{table}[H]
\centering
\begin{tabularx}{\textwidth}{lX}
\toprule
\textbf{Juego} & \textbf{Formato} \\
\midrule
FC 25 & Single-elimination: el que pierde queda eliminado \\
Call of Duty Black Ops 2 & Double-elimination: el que pierde va al bracket de perdedores; se elimina al perder dos veces \\
\bottomrule
\end{tabularx}
\caption{Formatos de bracket por juego}
\end{table}

\section{Estructura de archivos generada}

Al completar el flujo, la estructura del proyecto result\'o as\'i:

\begin{lstlisting}
sistema-torneos/
|-- server/
|   |-- src/
|   |   |-- controllers/
|   |   |   |-- auth.controller.ts
|   |   |   |-- inscripciones.controller.ts
|   |   |   |-- admin-inscripciones.controller.ts
|   |   |   `-- admin-bracket.controller.ts
|   |   |-- routes/
|   |   |-- services/
|   |   `-- middleware/
|   `-- prisma/
|       |-- schema.prisma
|       `-- migrations/      <- 3 migraciones aplicadas
|-- client/
|   `-- src/
|       |-- components/
|       |   |-- admin/
|       |   |-- brackets/
|       |   `-- inscripcion/
|       |-- pages/
|       `-- services/
|-- shared/
|   `-- types/
|       |-- auth.ts
|       |-- bracket.ts
|       `-- inscripcion.ts
`-- specs/
    |-- 001-inscripcion-torneo/
    |-- 002-login-admin/
    |-- 003-panel-admin-jugadores/
    |-- 004-generar-bracket-torneo/
    `-- 005-registrar-resultados-bracket/
\end{lstlisting}

Cada carpeta dentro de \texttt{specs/} contiene: \texttt{spec.md},
\texttt{plan.md}, \texttt{tasks.md}, \texttt{research.md},
\texttt{data-model.md} y una subcarpeta \texttt{contracts/}.

% ════════════════════════════════════════════════════════════════════════════
\chapter{Gesti\'on de Cambios en SDD}

\section{Agregar una feature nueva}

Cualquier funcionalidad nueva debe pasar por el flujo correspondiente
antes de generar c\'odigo:

\begin{figure}[H]
\centering
\begin{tikzpicture}[
  nodo/.style={
    rectangle, rounded corners=3pt,
    minimum width=2.6cm, minimum height=0.85cm,
    draw=gris, fill=grisCaja,
    font=\small\ttfamily, align=center
  },
  flecha/.style={-Stealth, thick, color=gris}
]
\node[nodo] (s) {specify};
\node[nodo, right=1cm of s] (p) {plan};
\node[nodo, right=1cm of p] (t) {tasks};
\node[nodo, right=1cm of t] (i) {implement};
\draw[flecha] (s)--(p);
\draw[flecha] (p)--(t);
\draw[flecha] (t)--(i);
\end{tikzpicture}
\caption{Flujo completo para una feature nueva}
\end{figure}

\section{Pasos seg\'un el tipo de cambio}

\begin{table}[H]
\centering
\begin{tabularx}{\textwidth}{lX}
\toprule
\textbf{Tipo de cambio} & \textbf{Pasos necesarios} \\
\midrule
Feature nueva compleja   & specify $\to$ plan $\to$ tasks $\to$ implement \\
Feature nueva simple     & specify $\to$ tasks $\to$ implement \\
Modificaci\'on peque\~na & specify $\to$ implement \\
Correcci\'on de bug      & implement (describiendo el bug con precisi\'on) \\
\bottomrule
\end{tabularx}
\caption{Pasos requeridos seg\'un el tipo de cambio}
\end{table}

\begin{notacaja}
  \textbf{Regla fundamental:} Nunca se modifica el c\'odigo directamente
  para agregar algo nuevo. El cambio nace en la spec, no en el c\'odigo.
  Esta regla es lo que diferencia SDD del vibe coding.
\end{notacaja}

% ════════════════════════════════════════════════════════════════════════════
\chapter{Alternativas a Spec Kit}

\section{Panorama de herramientas SDD en 2026}

\begin{table}[H]
\centering
\begin{tabularx}{\textwidth}{llX}
\toprule
\textbf{Herramienta} & \textbf{Complejidad} & \textbf{Mejor para} \\
\midrule
GitHub Spec Kit & Media & Aprendizaje, proyectos medianos, greenfield \\
GSD             & Baja  & Solo dev, flujo diario r\'apido, Claude Code \\
BMAD Method     & Alta  & Equipos grandes, proyectos enterprise \\
OpenSpec        & Baja  & Proyectos existentes (brownfield) \\
Amazon Kiro     & Media & Stacks AWS, IDE integrado con tier gratuito \\
\bottomrule
\end{tabularx}
\caption{Comparaci\'on de herramientas SDD principales en 2026}
\end{table}

\section{Spec Kit vs.\ GSD}

\begin{table}[H]
\centering
\begin{tabularx}{\textwidth}{lXX}
\toprule
\textbf{Aspecto} & \textbf{Spec Kit} & \textbf{GSD} \\
\midrule
Ceremonia    & Alta --- m\'ultiples archivos por feature & Baja --- minimalista \\
Velocidad    & M\'as lento, m\'as riguroso & M\'as r\'apido, m\'as pragm\'atico \\
Archivos     & spec, plan, tasks, research por feature & Contexto en un flujo \'unico \\
Ideal para   & Aprender SDD, proyectos complejos & Flujo diario, solo dev \\
Integraci\'on & 30+ agentes & Claude Code, Gemini CLI \\
\bottomrule
\end{tabularx}
\caption{Comparaci\'on directa: Spec Kit vs.\ GSD}
\end{table}

% ════════════════════════════════════════════════════════════════════════════
\chapter{Conclusiones y Pr\'oximos Pasos}

\section{Aprendizajes obtenidos}

\begin{enumerate}[leftmargin=*, label=\arabic{*}.]
  \item \textbf{SDD invierte la relaci\'on con el agente.} El desarrollador
        define qu\'e construir; el agente ejecuta bajo esas restricciones.
  \item \textbf{El \texttt{constitution.md} es el artefacto m\'as valioso.}
        Define el ADN del proyecto y persiste a lo largo de toda la
        implementaci\'on.
  \item \textbf{Cada fase requiere revisi\'on activa.} El flujo no es
        autom\'atico. Aprobar sin leer equivale a vibe coding disfrazado.
  \item \textbf{La spec por feature es la unidad de trabajo.} No se codea
        sin spec. No se avanza al plan sin aprobar la spec.
  \item \textbf{Los cambios nacen en la spec, no en el c\'odigo.}
        Modificar el c\'odigo directamente rompe la trazabilidad.
\end{enumerate}

\section{Pr\'oximo paso: GSD}

El siguiente paso es \textbf{GSD (Get Shit Done)}, el framework de
meta-prompting y context engineering construido principalmente para
Claude Code. GSD representa la versi\'on liviana del mismo principio SDD,
orientada al flujo diario de un desarrollador individual.

GSD ser\'a la herramienta utilizada en el desarrollo de \textbf{SABERHUB},
el LMS open source en construcci\'on, dado que su bajo overhead ceremonial
lo hace m\'as adecuado para iteraciones frecuentes en un proyecto de largo
plazo.

\vspace{1cm}

\begin{notacaja}
  \textbf{Conclusi\'on general:} El vibe coding resuelve el problema de
  velocidad inicial a costa de la calidad sostenida. SDD, implementado
  con Spec Kit o GSD, mantiene la velocidad en el largo plazo porque el
  agente siempre tiene contexto preciso. La inversi\'on en la fase de
  especificaci\'on se recupera ampliamente en la fase de implementaci\'on.
\end{notacaja}

\end{document}